% ============================================================================
% Hypothesis H2: spatial aggregation bias in robust site selection
% Pre-registered decision rule (OQ-050, H2), fixed 2026-10-13, before any
% real-country result was produced. Primary threshold gamma = 0.30.
% Notation follows docs/METHODOLOGY.md and the H1 section (sec:h1).
% Citation keys marked "verify" must be checked before submission.
% ============================================================================

\section{Hypothesis H2: Spatial Aggregation Bias in Robust Site Selection}
\label{sec:h2}

\subsection{Purpose and Position in the Research Design}
\label{subsec:h2_purpose}

Hypothesis H2 asks whether assessing renewable expansion at the level of
administrative regions, instead of at the level of individual cells, leads to
the selection of different and less robust locations. The question is not
whether resources and costs vary within a country, which is true of any large
territory, but whether that variation is large enough to change a siting
decision. H2 is therefore tested through a decision outcome: two procedures
select the same amount of capacity, one ranking cells and one ranking regions,
and the agreement between their selections is measured.

The question is grounded in the Modifiable Areal Unit Problem (MAUP), the
dependence of spatial results on the scale and the boundaries of the units used
to aggregate observations \cite{openshaw1984,gotwayyoung2002}. It is related to,
but distinct from, the ecological fallacy, which concerns inferring
individual-level properties from aggregate relationships \cite{robinson1950}:
H2 does not infer anything about cells from regional statistics, but compares
the selections that each level of representation would produce. In energy
system modelling, spatial resolution has been shown to affect the placement of
renewable generation, the representation of network constraints and the
estimated cost of expansion \cite{hoerschbrown2017,frysztacki2021,trondle2020},
and spatial representation is identified as a central methodological choice in
reviews of energy system models \cite{pfenninger2014,javanmardi2025}.
% verify javanmardi2025

The mechanism tested is simple. A regional score summarises many cells. A
region whose average is mediocre may contain the best sites of the country,
diluted by poor ones; a region with a good average may contain no excellent
site at all. Ranking whole regions by their average can therefore select
capacity that a cell-level ranking would not choose. H2 measures how much of
the cell-level selection survives regional aggregation.

\subsection{Notation}
\label{subsec:h2_notation}

The notation of Section~\ref{sec:h1} applies. For one country and one
technology, $\mathcal{I}$ is the set of cells entering the robustness analysis,
with $N=|\mathcal{I}|$. For each cell $i\in\mathcal{I}$:
\begin{itemize}
    \item $q_i$ is the installable capacity (MW) of the central land scenario,
    computed by GeoFREA from the eligible area, the land-utilisation factor and
    the power density; it does not depend on the climate member;
    \item $z_i=\mathrm{MR}_i$ is the primary robustness score of the thesis,
    the maximum over climate members of the core window of the 90th percentile
    over parameter samples of the relative regret
    (Equation~\ref{eq:h1_relative_regret}). Lower values are better.
\end{itemize}
The partition $\mathcal{A}=\{A_1,\dots,A_M\}$ groups the cells of
$\mathcal{I}$ by first-level administrative unit (Admin1: states in Brazil and
India, districts in mainland Portugal), using the Admin1 identifier assigned to
each cell by GeoFREA. Regions without cells in $\mathcal{I}$ do not enter the
analysis. The capacity of region $A_m$ is $Q_m=\sum_{i\in A_m}q_i$ and the total
candidate capacity is $Q_{\max}=\sum_{i\in\mathcal{I}}q_i$.

Ranks are assigned in ascending order of score with the deterministic tie
convention of the F7 phase record, applied identically to cells and regions.

\subsubsection{Capacity target}

The target capacity is the capacity of the robust top-$k$ set used in H1:
\begin{equation}
    T=\sum_{i\in S_k}q_i,
    \qquad
    S_k=\{\,k \text{ cells of } \mathcal{I} \text{ with the lowest } z_i\,\},
    \qquad
    k=\lceil p_k N\rceil .
    \label{eq:h2_target}
\end{equation}
Defining $T$ through the same top-$k$ set links H1 and H2 to one shortlist and
guarantees that the target exists for every country and technology. National
capacity targets are used as a sensitivity where they are available.

\subsection{The Two Selections}
\label{subsec:h2_selections}

\subsubsection{Cell-level selection}

Cells are taken in ascending order of $z_i$ until the cumulative capacity
reaches $T$. By construction of Equation~\ref{eq:h2_target}, this selection is
exactly $S_k$. Its selection vector is $x^{\mathrm{cell}}_i=1$ for $i\in S_k$
and $0$ otherwise, so that $\sum_i q_i x^{\mathrm{cell}}_i=T$.

\subsubsection{Regional selection}

Each region receives the capacity-weighted mean score
\begin{equation}
    \overline{z}_m=\frac{\sum_{i\in A_m}q_i z_i}{Q_m}.
    \label{eq:h2_region_score}
\end{equation}
Capacity weighting is the summary a regional study would naturally use: a
region's score reflects where its capacity is. Regions are taken whole, in
ascending order of $\overline{z}_m$, until the target is reached. The last
region $A_{m^{*}}$ enters with the fraction
\begin{equation}
    y_{m^{*}}=\frac{T-\sum_{m\in\mathcal{M}_{\mathrm{prev}}}Q_m}{Q_{m^{*}}},
    \qquad 0<y_{m^{*}}\leq 1,
    \label{eq:h2_fractional_region}
\end{equation}
where $\mathcal{M}_{\mathrm{prev}}$ is the set of regions taken before it. Every
cell inherits the fraction of its region, $x^{\mathrm{admin}}_i=y_m$ for
$i\in A_m$, with $y_m=1$ for the regions taken whole and $y_m=0$ for those not
taken, so that $\sum_i q_i x^{\mathrm{admin}}_i=T$.

The fractional last region is an accounting device that makes the two
selections exactly equal in capacity. It does not represent how projects would
be distributed inside that region, and it applies the same fraction to all of
its cells.

\subsection{Primary Statistic: Capacity Overlap Index}
\label{subsec:h2_overlap}

The capacity common to both selections is
\begin{equation}
    \mathrm{Cap}_{\mathrm{common}}
    =\sum_{i\in\mathcal{I}}q_i\min\!\left(x^{\mathrm{cell}}_i,x^{\mathrm{admin}}_i\right),
    \label{eq:h2_common_capacity}
\end{equation}
and the Capacity Overlap Index is
\begin{equation}
    O=\frac{\mathrm{Cap}_{\mathrm{common}}}{T},\qquad 0\leq O\leq 1 .
    \label{eq:h2_overlap}
\end{equation}
$O$ is the share of the robust cell-level capacity that the regional procedure
also selects. Because both selections total $T$, $1-O$ is at the same time the
share of the cell-level selection missed by the regional one and the share of
the regional selection that the cell-level ranking would not choose.

\subsubsection{Reference value}

If the regional procedure carried no information about where the robust cells
are, it would behave like a single national region (Admin0): every cell would
receive the same fraction $T/Q_{\max}$, and the overlap would be
\begin{equation}
    O_{\mathrm{ref}}=O_{\mathrm{Admin0}}=\frac{T}{Q_{\max}} .
    \label{eq:h2_reference}
\end{equation}
$O_{\mathrm{ref}}$ is the agreement obtained by spreading the target uniformly
over the country. It is reported with every result, and it defines when the
test is informative (Section~\ref{subsec:h2_rule}).

\subsection{Hypothesis and Decision Rule}
\label{subsec:h2_rule}

H2 states that the Admin1 selection shares little of the capacity of the
spatially explicit robust selection. The verdict for each country and
technology is:
\begin{equation}
    \text{H2 is}\;
    \begin{cases}
        \text{not testable}, & O_{\mathrm{ref}}\geq\gamma,\\
        \text{supported}, & O_{\mathrm{ref}}<\gamma \text{ and } O<\gamma,\\
        \text{not supported}, & O_{\mathrm{ref}}<\gamma \text{ and } O\geq\gamma,
    \end{cases}
    \qquad \gamma=0.30 .
    \label{eq:h2_rule}
\end{equation}

The threshold $\gamma=0.30$ requires that less than three-tenths of the robust
cell-level capacity be selected by the regional procedure; equivalently, more
than seven-tenths of it is lost to aggregation. It is deliberately stricter
than a majority criterion ($\gamma=0.50$): a lower threshold makes it harder to
conclude that aggregation distorts the selection, so a supported verdict
requires strong disagreement between the two procedures. The strict inequality
means that $O=0.30$ does not support H2. The rule is deterministic, applied to a
computed index; it is not a statistical test with a sampling distribution, and
its stringency is a property of the threshold, not a reduction of a Type I error
rate.

The not-testable case protects the rule from a degenerate outcome. When the
target is a large share of the candidate capacity, even a uniform spread over
the country reaches an overlap of $\gamma$ or more, and a low overlap can no
longer be distinguished from the effect of the target size. For the top-$k$
shares used in the thesis $O_{\mathrm{ref}}$ is expected to lie well below
$0.30$, but the condition is checked and reported for every case.

The value of $\gamma$ is a decision threshold of this thesis. The MAUP and
spatial-resolution literature establishes that aggregation can change results,
but provides no universal overlap threshold. The defence of $\gamma=0.30$ rests
on its being fixed before any real result was produced, on its direct reading
(less than three-tenths shared), and on the decision profile reported in
Section~\ref{subsec:h2_sensitivity}.

\subsection{Complementary Diagnostics (No Verdict)}
\label{subsec:h2_diagnostics}

Three diagnostics describe the nature of the mismatch. None of them changes the
verdict.

\subsubsection{Robustness penalty of the regional selection}

\begin{equation}
    \Delta_z=\frac{1}{T}\sum_{i\in\mathcal{I}}q_i x^{\mathrm{admin}}_i z_i
    -\frac{1}{T}\sum_{i\in\mathcal{I}}q_i x^{\mathrm{cell}}_i z_i .
    \label{eq:h2_penalty}
\end{equation}
$\Delta_z$ is the increase in the capacity-weighted maximum relative regret
caused by selecting regions instead of cells, in the same units as
$\mathrm{MR}$. Since the cell-level selection takes the lowest scores first,
$\Delta_z\geq 0$ up to ties. A low overlap with a small $\Delta_z$ indicates
that the two procedures choose different but similarly robust capacity; a low
overlap with a large $\Delta_z$ indicates a material loss of robustness.

\subsubsection{Within-region heterogeneity}

\begin{equation}
    \sigma^{2}_{z,m}=\frac{\sum_{i\in A_m}q_i\left(z_i-\overline{z}_m\right)^{2}}{Q_m}.
    \label{eq:h2_within_variance}
\end{equation}
Large within-region dispersion is the mechanism through which aggregation can
hide the best sites, but it does not by itself imply a low overlap, which also
depends on the ordering of the regional means and on where the capacity lies.
The interquartile range of $z_i$ over the country, relative to its median, is
reported alongside.

\subsubsection{Regional distribution of the selected capacity}

With $P^{\mathrm{cell}}_m$ and $P^{\mathrm{admin}}_m$ the shares of $T$ selected
in region $A_m$ by each procedure, the total variation distance is
\begin{equation}
    D_{\mathrm{TV}}=\frac{1}{2}\sum_{m=1}^{M}
    \left|P^{\mathrm{cell}}_m-P^{\mathrm{admin}}_m\right| .
    \label{eq:h2_tv}
\end{equation}
It shows whether the mismatch is between regions (high $D_{\mathrm{TV}}$) or
within the same regions (low $D_{\mathrm{TV}}$ with low $O$).

\subsection{Sensitivity Analyses}
\label{subsec:h2_sensitivity}

The verdict is issued from the primary configuration only. A change of verdict
in a variant is reported as a dependence of the result on that choice, not as a
revised verdict.

\begin{enumerate}
    \item Decision profile. Since $O$ does not depend on $\gamma$, the profile
    reports the verdict for $\gamma\in\{0.10,0.15,\dots,0.70\}$, with $0.50$
    marked as the majority benchmark. H2 is supported for every
    $\gamma>O$, so the profile shows the smallest threshold that would still
    support it.
    \item Target size. $T$ is recomputed for $p_k\in\{0.05,0.10,0.20\}$ and, where
    an official national target is available for the technology and horizon,
    for that target. The overlap converges to 1 as $T$ approaches $Q_{\max}$;
    no other monotonic behaviour is assumed.
    \item Regional ranking statistic. Regions ranked by the capacity-weighted
    median of $z_i$ instead of the mean.
    \item Climate configuration. $z_i$ recomputed without IPSL-CM6A-LR, the
    climate model retained as a recorded exception, and for the 2071--2100
    window.
    \item Cell scale. The analysis repeated at 0.1 degree, which separates the
    effect of the cell resolution from that of the regional aggregation.
\end{enumerate}
Finer administrative partitions (Admin2) and electricity market zones are not
part of the present design, because they are not available consistently for the
three countries; their absence is stated as a limitation.

\subsection{Worked Example}
\label{subsec:h2_example}

A country has $Q_{\max}=1{,}000$~MW of candidate capacity in three states, and
the robust top-$k$ set holds $T=100$~MW, so $O_{\mathrm{ref}}=0.10<0.30$ and the
test is informative.

\begin{itemize}
    \item State A holds 60~MW of the most robust cells and 300~MW of poor ones;
    its capacity-weighted mean is mediocre.
    \item State B holds 400~MW of uniformly intermediate cells and no cell of
    the top-$k$; its mean is the best of the three.
    \item State C holds 40~MW of top-$k$ cells and 200~MW of poor ones.
\end{itemize}
The cell-level selection takes the 60~MW of A and the 40~MW of C. The regional
selection ranks B first and takes 100~MW from it, a fraction $y_B=0.25$ of its
capacity. No selected capacity is common, so $O=0$ and H2 is supported: the
regional averages of A and C dilute their best sites, and B, without a single
top site, wins on its average. $\Delta_z$ then measures how much robustness is
lost by building in B.

If instead 42~MW were common, $O=0.42$: H2 would not be supported under
$\gamma=0.30$, although it would be under the majority benchmark $0.50$. The
decision profile makes this dependence explicit. The values are illustrative
and are not GeoFREA results.

\subsection{Parameters}
\label{subsec:h2_parameters}

\begin{table}[htbp]
    \centering
    \caption{Pre-registered parameters of H2.}
    \label{tab:h2_parameters}
    \small
    \renewcommand{\arraystretch}{1.2}
    \begin{tabularx}{\textwidth}{
        @{}
        >{\raggedright\arraybackslash}p{2.4cm}
        >{\raggedright\arraybackslash}p{3.0cm}
        >{\raggedright\arraybackslash}X
        @{}
    }
        \toprule
        \textbf{Parameter} & \textbf{Value} & \textbf{Role} \\
        \midrule
        $\gamma$ & 0.30 & Primary threshold; H2 is supported when less than
        three-tenths of the robust cell-level capacity is selected by the
        regional procedure. \\
        \addlinespace
        $z_i$ & $\mathrm{MR}_i$ & Primary robustness score; lower is better. \\
        \addlinespace
        $q_i$ & Central-scenario capacity & Capacity weight of each cell. \\
        \addlinespace
        $T$ & Capacity of the top-$k$ set & Same shortlist as H1;
        $k=\lceil p_k N\rceil$, $p_k$ fixed in OQ-021. \\
        \addlinespace
        $\mathcal{A}$ & Admin1 & Primary regional partition. \\
        \addlinespace
        $\overline{z}_m$ & Capacity-weighted mean & Regional ranking
        statistic. \\
        \addlinespace
        $O_{\mathrm{ref}}$ & $T/Q_{\max}$ & Overlap of a uniform national
        spread; the test is not testable when $O_{\mathrm{ref}}\geq\gamma$. \\
        \bottomrule
    \end{tabularx}
\end{table}

\subsection{Limitations}
\label{subsec:h2_limitations}

Dependence on the target. $O$ is conditional on $T$; the sensitivity to the
top-$k$ share and to national targets shows whether the conclusion holds across
expansion scales.

Dependence on the score. The regional and cell-level selections are both
determined by $\mathrm{MR}_i$. H2 evaluates the consequence of aggregation for
this score; it does not establish that the score is the only valid measure of
robustness. The cell-level selection is the benchmark of the GeoFREA framework,
not an independently verified optimum.

Meaning of a low overlap. A low $O$ shows that the two procedures choose
different capacity. Whether this implies higher system costs, transmission needs
or curtailment requires the robustness penalty $\Delta_z$ and, beyond this
thesis, system-level modelling.

Few regions. Countries with few Admin1 units, such as mainland Portugal, give a
coarse regional selection, and the result there depends more on the boundaries
than on the scale. The reference value and the regional distribution distance
are read together with $O$ in those cases.

Partitions. Only the Admin1 partition is tested. Administrative units are not
necessarily the units of energy planning, and other partitions at similar scale
may give different overlaps (zoning effect).

Fractional accounting. The fractional last region equalises capacity for the
comparison and does not represent indivisible projects, minimum plant sizes or
network constraints.

Generalisation. Results are conditional on the country, technology, candidate
set, target and partition, and are reported per country and technology.
